\section{Éléments contextuels}

\begin{guideline}
Cette section doit être concise (idéalement 1 page recto-verso max). Elle permet aux évaluateurs d'appréhender rapidement le contexte de votre stage. Soyez factuel·le et précis·e. Pensez à produire également un executive summary en anglais.
\end{guideline}


\subsection{Présentation de l'entreprise}

\begin{guideline}
Décrivez l'entreprise en quelques lignes : nom, secteur d'activité, taille (effectifs, CA), type d'organisation (startup, grand groupe, PME...), positionnement marché. L'objectif est de donner une image claire en un coup d'œil, pas de rédiger une brochure commerciale.
\end{guideline}

\begin{center}
\renewcommand{\arraystretch}{1.5}
\begin{tabularx}{\textwidth}{@{}l X@{}}
\textbf{Nom} & \instruction{Ex. : Accenture France} \\
\textbf{Secteur} & \instruction{Ex. : Conseil en transformation digitale} \\
\textbf{Taille} & \instruction{Ex. : \textasciitilde500 000 employés dans le monde, \textasciitilde10 000 en France} \\
\textbf{Type} & \instruction{Ex. : Multinationale cotée / Startup SaaS / ETI familiale...} \\
\textbf{Lieu du stage} & \instruction{Ex. : Paris La Défense (siège) --- télétravail partiel} \\
\end{tabularx}
\end{center}

\instruction{[Complétez ici avec une brève présentation narrative de l'entreprise, son positionnement, ses produits/services clés --- 5 à 10 lignes]}


\subsection{L'équipe intégrée}

\begin{guideline}
Décrivez l'équipe dans laquelle vous avez travaillé : taille, composition, rôles des membres, votre positionnement au sein de l'équipe. Vous pouvez inclure un organigramme simplifié si cela apporte de la clarté.
\end{guideline}

\instruction{[Décrivez l'équipe : effectif, organisation, votre rôle et votre positionnement. Ex. : « J'ai intégré une équipe de 8 personnes composée de 3 ingénieurs, 2 data scientists, 1 chef de projet et 1 product owner. J'étais rattaché·e directement au lead data scientist. »]}


\subsection{Mission de stage}

\begin{guideline}
Décrivez la mission initiale telle qu'elle vous a été confiée au démarrage, puis listez toutes les évolutions qui ont eu lieu (ajout, modification, recentrage...) avec leurs dates. Cela montre votre capacité à vous adapter.
\end{guideline}

\begin{center}
\renewcommand{\arraystretch}{1.5}
\begin{tabularx}{\textwidth}{@{}l X@{}}
\textbf{Mission initiale} & \instruction{[Décrivez la mission de départ telle que définie dans la convention de stage]} \\
\textbf{Évolution(s) datée(s)} & \instruction{[JJ/MM] --- [Description de l'évolution]} \\
\end{tabularx}
\end{center}
